\documentclass[11pt]{article}
\usepackage[margin=1.1in]{geometry}
\usepackage{amsmath}
\usepackage{booktabs}
\usepackage{microtype}
\usepackage[colorlinks=true,allcolors=blue]{hyperref}

\title{\textbf{The Empirical Case for Gaussian Renormalization as Default}}
\author{Peter Cotton}
\date{Provisional draft, 17 August 2026}

\begin{document}
\maketitle

\begin{abstract}
\noindent A population's choice shares on a full menu do not tell you its choices
from smaller menus without a rule connecting them. Luce's axiom supplies one:
renormalize, leaving the odds between survivors untouched. An independent Gaussian
race supplies another: re-run the contest without the removed options, which makes
the survivors less lopsided. Neither rule needs data from the smaller menus. In an
experiment where subjects chose separately from all thirty-one subsets of five goods,
the Gaussian rule predicts held-out choices better, and across twelve archival
ranking collections it does so in twelve of twelve. Both results survive comparison
with a null in which the axiom holds by construction, which matters because
contraction of a noisily estimated share vector is also shrinkage. The recommendation
is narrow and practical: when mapping aggregate full-menu shares to restricted-menu
shares, make the Gaussian race the default.
\end{abstract}

\section{Two rules, one calibration}

For a population and a menu $S$, let $p_i$ be the share choosing or first-ranking
item $i$. Luce's axiom \cite{luce1959} assigns a worth $u(i)$ with
\begin{equation}
p_i = \frac{u(i)}{\sum_{j \in S} u(j)},
\label{eq:luce}
\end{equation}
so the shares determine the worths up to scale, and the prediction for any subset
$T \subseteq S$ follows by changing the denominator alone. Odds between survivors are
untouched: for $i,j \in T$ the share preferring $i$ is $p_i/(p_i+p_j)$ when $T$ is the
pair. This is Independence of Irrelevant Alternatives.

Thurstone \cite{thurstone1927} connects the same two quantities through a contest.
Item $i$ carries a location $a_i$, draws a performance $X_i \sim N(a_i,1)$ on each
occasion, and the highest draw is chosen, so
\begin{equation}
p_i = \int \varphi(x - a_i) \prod_{k \neq i} \Phi(x - a_k)\, dx.
\label{eq:winprob}
\end{equation}
The integral is one-dimensional for any number of alternatives, since conditioning on
$X_i = x$ makes the other events independent, and it is inverted numerically by a
lattice method \cite{cotton2021}; locations are identified up to a common translation,
fixed here by $\sum_i a_i = 0$. Prediction for a subset $T$ means re-running the
contest among $T$ alone. For a pair this is one normal evaluation,
\begin{equation}
P(X_i > X_j) = \Phi\!\left(\frac{a_i - a_j}{\sqrt{2}}\right),
\label{eq:pair}
\end{equation}
and it is less lopsided than $p_i/(p_i+p_j)$: winning a large field needs an unusually
good draw, and a favourite's edge is partly an edge in producing extremes, so that
edge counts for less when one rival remains.

Both rules are calibrated to the same $|S|-1$ free shares and neither sees any
observation from a smaller menu. The comparison is therefore between two
zero-extra-parameter maps, not between two fitted models. That is worth stating
plainly, because it also bounds what the comparison can show. A fitted
Plackett--Luce likelihood uses the whole ranking and a fitted Bradley--Terry model
uses the pairwise outcomes directly; neither is tested here. What is tested is a
cross-moment implication of the homogeneous Luce model, namely whether its full-menu
shares correctly determine its restricted-menu shares.

That the direction of the Gaussian disagreement is not an accident can be shown. With
the highest draw winning, let $r = f/F$ be the reverse hazard of the common noise. For
nested menus $T \subset S$ and $i,j \in T$ with $a_i > a_j$, if $\log r$ is concave
then the odds of $i$ against $j$ are at least as large in $S$ as in $T$, strictly so
when $\log r$ is strictly concave and at least one contestant is actually removed. The
argument is a covariance inequality between two increasing functions, and the Gaussian
satisfies the hypothesis because
$\frac{d^2}{dx^2}\log r(x) = -\operatorname{Var}(Z \mid Z < x) < 0$. The Gumbel has
$r(x) = s^{-1}e^{-x/s}$, so $\log r$ is affine, the inequality is an equality, and
renormalization is recovered exactly; this is the sufficiency half of Yellott's
theorem \cite{yellott1977}, which is what makes the axiom a boundary of the Gaussian's
class rather than a rival to it. We use the proposition only for the sign. It gives no
ordering of magnitudes, and contraction is not a property of a noise law at all: it
depends on the share configuration, and laws recalibrated to the same shares reorder
between configurations, so we make no claim of the form heavier tails contract less.

\section{What is observed, and at what level}
\label{sec:estimand}

Twelve of the thirteen collections used here are complete rankings. A ranking yields
both the first-placed item and the relative order of any two items, so in those twelve
the comparison is between two accounts of one ranking distribution rather than a
record of behaviour under restriction. A respondent's ranking imposes one ordering on
every subset by construction, so no ranking dataset can show a person choosing $i$
over $j$ from one menu and $j$ over $i$ from another, and none reveals whether that
person would choose as their ranking implies when handed two options. Any statement
about restriction from those twelve rests on the assumption that they would.

The thirteenth makes no such assumption, and for that reason it comes first below.

The level of analysis matters as much as the elicitation. Every quantity here is a
population quantity, and aggregate failure of the axiom does not imply that
individuals violate it. A population of people who each choose by
Equation~\eqref{eq:luce} with worths differing between people generally does not
satisfy the axiom in aggregate, which is why mixed logit is useful and why its
probabilities are integrals over logit probabilities rather than logit probabilities
\cite{mcfaddentrain2000}. So what can be rejected here is a representative-agent Luce
restriction map. Individual-level noise is not identified.

\section{Choices from real menus}
\label{sec:menus}

Costa-Gomes, Cueva, Gerasimou and Tejiscak \cite{costagomes2022} had each subject
choose separately from all thirty-one subsets of five consumer goods, so restriction is
observed rather than imputed. The archive names the goods only by two-letter codes and
we describe them no more precisely. This is the cleanest available test of the question
above, and it is the primary result.

Subjects are split into five folds. On the training subjects, shares come only from the
full five-item menu; both rules are calibrated to those shares and nothing else; each
then predicts every menu of size two, three and four; and we score the choices
actually made by held-out subjects. Intervals resample subjects and repeat the whole
procedure, calibration included. Some subjects may defer, so the forced-choice
subgroup is reported separately: those subjects cannot defer, and no coding decision
about deferral enters that row.

\begin{table}[ht]
\centering
\small
\begin{tabular}{lrrrrl}
\toprule
Subjects & $n$ & Choices & Renorm. & Race & Gain \\
\midrule
All & 364 & 8{,}239 & $0.9263$ & $0.9164$ & $+0.0100$ $[+0.0048, +0.0246]$ \\
Forced choice only & 137 & 3{,}425 & $0.9112$ & $0.8846$ & $+0.0265$ $[+0.0116, +0.0700]$ \\
\bottomrule
\end{tabular}
\caption{Held-out log loss per choice on menus of size two to four, with both rules
calibrated only to training-fold full-menu shares under a common add-$\alpha$
convention, $\alpha = 1/2$. Shares rest on $268$ full-menu choices for all subjects and
$109$ for the forced-choice subgroup. Intervals are subject bootstraps that recompute
folds, shares, calibration and predictions.}
\label{tab:menus}
\end{table}

The Gaussian rule is better on both, and better by more where deferral cannot
interfere. Section~\ref{sec:null} shows the gain is not an artefact of estimating
shares from a few hundred observations.

\section{Choices induced from rankings}
\label{sec:rankings}

Twelve archival collections have human respondents and rankings, or ratings fine enough
to yield them. Two rank-order card tasks from the General Social Survey ask what a
child should learn and what matters in a job. Three independent samples of Jester users
rate the same ten jokes. Sushi rankings come from a Japanese web survey, film rankings
are induced from Netflix ratings, and dots and puzzles are perceptual discrimination
tasks. The rest are rankings of seven sports by participation preference, ten
occupations by perceived prestige, and four political goals from a five-nation survey.
The thirteen collections comprise $147{,}719$ responses, of which the menu experiment
contributes $373$ subjects; the twelve ranking collections contribute $147{,}346$.
These are not thirteen independent experiments, since three are Jester files, two are
General Social Survey items, and dots and puzzles are companion tasks.

The protocol matches Section~\ref{sec:menus} with respondents in place of subjects, and
the outcome for a subset is the respondent's highest-ranked member of it.

\begin{table}[ht]
\centering
\small
\begin{tabular}{lrrrrrl}
\toprule
Dataset & $n$ & $K$ & Subsets & Renorm. & Race & Gain \\
\midrule
Occupational prestige & 143 & 10 & 1{,}013 & $1.0173$ & $0.9783$ & $+0.0390$ $[+0.0227, +0.0555]$ \\
Sushi & 5{,}000 & 10 & 1{,}013 & $1.3724$ & $1.3613$ & $+0.0111$ $[+0.0094, +0.0128]$ \\
Political goals & 2{,}262 & 4 & 11 & $0.8835$ & $0.8766$ & $+0.0069$ $[+0.0059, +0.0078]$ \\
GSS socialization & 5{,}000 & 5 & 26 & $0.7984$ & $0.7929$ & $+0.0055$ $[+0.0041, +0.0070]$ \\
Jester file 3 & 5{,}000 & 10 & 1{,}013 & $1.5189$ & $1.5165$ & $+0.0024$ $[+0.0015, +0.0032]$ \\
Jester file 1 & 5{,}000 & 10 & 1{,}013 & $1.5287$ & $1.5267$ & $+0.0020$ $[+0.0012, +0.0027]$ \\
Jester file 2 & 5{,}000 & 10 & 1{,}013 & $1.5247$ & $1.5230$ & $+0.0018$ $[+0.0010, +0.0025]$ \\
GSS job values & 5{,}000 & 5 & 26 & $0.8593$ & $0.8576$ & $+0.0017$ $[+0.0005, +0.0030]$ \\
Dots & 3{,}183 & 4 & 11 & $0.8285$ & $0.8270$ & $+0.0015$ $[+0.0005, +0.0025]$ \\
Netflix & 4{,}379 & 3 & 4 & $0.7932$ & $0.7922$ & $+0.0009$ $[+0.0007, +0.0012]$ \\
Puzzles & 3{,}180 & 4 & 11 & $0.8180$ & $0.8173$ & $+0.0007$ $[-0.0003, +0.0017]$ \\
\midrule
Sports participation & 130 & 7 & 120 & $1.2579$ & $1.2528$ & $+0.0050$ $[+0.0017, +0.0087]$ \\
\bottomrule
\end{tabular}
\caption{Held-out log loss per prediction over every subset of size two or more.
Datasets are capped at five thousand respondents for runtime. The race is lower in
twelve of twelve and the interval excludes zero in eleven. Occupational prestige is
scorable here only because the shared smoothing removes a zero training cell. Sports
participation sits below the rule: its gain is real but Section~\ref{sec:null} shows it
carries no evidence about the restriction map. Intervals resample respondent losses
with the fitted training models held fixed, so unlike Table~\ref{tab:menus} they omit
calibration uncertainty and are too narrow, most seriously in the two smallest rows.}
\label{tab:heldout}
\end{table}

Averaging over every subset weights cardinalities by how many there are, so
intermediate menus dominate: for $K=10$ only $45$ of $1{,}013$ subsets are pairs, about
four per cent of the weight. Those middling menus are where least has been removed and
least is at stake. Splitting by menu size shows the pairwise gain exceeding the
aggregate in eleven of twelve datasets, by factors from $0.7$ to $3.7$ with median
$1.8$: Sushi is $+0.0412$ at $|T| = 2$ against $+0.0111$ aggregate, and GSS
socialization $+0.0130$ against $+0.0055$. Occupational prestige is the exception,
peaking at $|T| = 4$. At $|T| = K$ nothing has been removed, the two rules must agree,
and the measured gain is zero to four decimals in every dataset, which is the check
that the pipeline is sound. We keep the all-subsets aggregate as the headline because it
fixes the estimand in advance rather than selecting a cardinality after the fact.

\section{Is the gain structural?}
\label{sec:null}

A positive gain has two possible sources and log loss cannot tell them apart. The race
may be closer to how the population chooses. Or the race may be wrong about that and
win anyway, because contracting a noisily estimated share vector is shrinkage, and
shrinkage lowers log loss when the input is over-confident. The second mechanism needs
no departure from the axiom, so properness of the score licenses nothing here: a proper
score says the true probability vector is optimal, not that a plug-in estimate of it
beats a biased transformation of the same estimate.

The remedy is to build a population where the axiom holds by construction and run the
identical procedure on it. Taking the observed full-menu shares as worths, choices or
complete rankings are generated from that exact process, and everything is repeated.
Rankings are drawn by sorting $\log u_i + G_i$ for independent standard Gumbel $G_i$,
which is the Plackett--Luce ranking law; for $K \ge 4$ this is a particular joint law
consistent with those subset marginals rather than the only one, so the null is
specifically a Plackett--Luce ranking null.

\begin{table}[ht]
\centering
\small
\begin{tabular}{lrrrr}
\toprule
Dataset & Observed & Null median & Excess & $p$ \\
\midrule
Menu experiment, forced choice & $+0.0265$ & $-0.0041$ & $+0.0306$ & $0.010$ \\
Menu experiment, all subjects & $+0.0100$ & $-0.0023$ & $+0.0123$ & $0.015$ \\
\midrule
Occupational prestige & $+0.0390$ & $-0.0227$ & $+0.0617$ & $\le 0.005$ \\
Sushi & $+0.0111$ & $-0.0062$ & $+0.0174$ & $\le 0.005$ \\
GSS socialization & $+0.0055$ & $-0.0066$ & $+0.0121$ & $\le 0.005$ \\
Political goals & $+0.0069$ & $-0.0012$ & $+0.0081$ & $\le 0.005$ \\
GSS job values & $+0.0017$ & $-0.0047$ & $+0.0065$ & $\le 0.005$ \\
Jester file 3 & $+0.0024$ & $-0.0015$ & $+0.0039$ & $\le 0.005$ \\
Jester file 1 & $+0.0020$ & $-0.0014$ & $+0.0034$ & $\le 0.005$ \\
Jester file 2 & $+0.0018$ & $-0.0014$ & $+0.0032$ & $\le 0.005$ \\
Dots & $+0.0015$ & $-0.0014$ & $+0.0029$ & $\le 0.005$ \\
Puzzles & $+0.0007$ & $-0.0011$ & $+0.0019$ & $\le 0.005$ \\
Netflix & $+0.0009$ & $-0.0001$ & $+0.0010$ & $\le 0.005$ \\
\midrule
Sports participation & $+0.0050$ & $+0.0051$ & $-0.0000$ & $0.51$ \\
\bottomrule
\end{tabular}
\caption{The same gain when the axiom holds by construction. Two hundred replicates per
row, so $p$ is floored at $1/201$; the menu experiment uses its own $373$ subjects and
menus. Wherever shares rest on more than about two thousand observations the null is
negative, because contraction is then a bias rather than useful insurance, so comparing
an observed gain with zero understates it.}
\label{tab:null}
\end{table}

The shrinkage mechanism is real and it accounts for the entire gain in the smallest
dataset. Sports participation, with $130$ respondents, sits exactly on its null and is
excluded from every claim here. Everywhere else the mechanism runs against the race, and
the excess over the null exceeds the raw gain in all twelve remaining rows. It also
rescues puzzles, whose raw interval in Table~\ref{tab:heldout} covers zero while its
excess over the null does not. The menu experiment clears the null on both readings,
which is the result the paper rests on, since it is the only one where restriction is
observed.

\section{How much of the available improvement is that?}

An absolute gain of a few thousandths of a nat is a valid predictive effect size but a
poor guide to how much was on offer, because both rules receive the same shares and the
shares carry nearly all the predictive content of a restricted menu. To scale it we also
score a smoothed subset-frequency forecast, the training-fold frequency with which each
item is top within each subset, which is $|T|-1$ free probabilities per subset.

This benchmark is not an oracle ceiling, and treating it as one would overstate the
ratio. It is itself estimated, so its expected held-out loss exceeds the oracle
conditional entropy by roughly $(m-1)/2n$ for an $m$-outcome menu and $n$ training
observations, and a regularized or structured forecast can beat it out of sample.
Applying that first-order correction, the Gaussian race closes a median $35$ per cent of
the empirical benchmark gap, from $16$ per cent for sports participation to $58$ per
cent for dots. Uncorrected the median is $39$ per cent. Either way the reading is that a
rule with nothing fitted to restricted menus recovers a third or more of what a
per-subset frequency estimator achieves, while renormalization recovers none of it,
being the unique odds-preserving map. Denominators as small as $0.002$ make these ratios
unstable, and they are offered as context for the absolute gains rather than as a
headline.

\section{What this does and does not establish}

The practical recommendation is narrow. When the task is to infer restricted-menu
population shares from full-menu population shares, the Gaussian race should be the
default in place of renormalization. It costs one extra one-dimensional integral, needs
no additional data, and was better on every collection examined here except the one too
small to discriminate. In the single experiment where restriction is actually observed
rather than induced from rankings, it is better by $+0.0265$ nats per choice among
subjects who cannot defer.

Three things do not follow. Nothing here compares fitted likelihoods: a fitted
Plackett--Luce model uses the whole ranking and a fitted Bradley--Terry model uses
pairwise outcomes directly, and neither was estimated, so no recommendation about
ranking or paired-comparison likelihoods is supported. The held-out score is a sum of
marginal scores for top-item-within-subset outcomes, proper for those marginals but not
for the full ranking distribution; with $K = 4$ the scored marginals carry at most
$17$ degrees of freedom against $23$ for a distribution over rankings, so distinct
ranking laws can agree on everything scored. And nothing identifies individual-level
noise, for the reason given in Section~\ref{sec:estimand}: heterogeneous Luce choosers
reproduce aggregate contraction, so the target rejected is a representative-agent map.

Nesting is also not the same as costlessness. The independent-noise contest class
contains the Gumbel point, so permitting the class excludes no Luce population. But
Table~\ref{tab:null} shows that choosing the Gaussian point specifically is a
statistical decision that loses when the axiom is true and the shares are well
estimated. The class costs nothing in generality; the default within it is a bet, and
the evidence here is that it is a good one for human population data.

What would settle more: scoring the exact Gaussian ordering law against
Plackett--Luce ranking probabilities. We do not do that. Pricing ordered outcomes by
removing the winner and re-running is a sequential rule, not the race's ordering law,
and against that law it misprices individual pairs and triples by up to a factor of
$3.4$, so it cannot stand in.

\section*{Data and code}

Sushi rankings come from Kamishima \cite{kamishima2003}; Jester ratings from the
Eigentaste project \cite{goldberg2001}; the General Social Survey cumulative file from
NORC; the perceptual, leisure, prestige, film and political-goal rankings via PrefLib
\cite{mattei2013} and the CRAN packages PLMIX and ConsRank; the menu-choice data from
the replication archive of \cite{costagomes2022}. All are public. Analysis code and
every input needed to reproduce the tables are under \texttt{research/human} in
\texttt{github.com/microprediction/winning}, the inputs as column extracts rather than
raw archives.

Two storage conventions must be told apart when reproducing this. BayesMallows,
ConsRank and PerMallows store ranks, so element $j$ holds the rank of alternative $j$,
while PLMIX stores orderings, so the first element names the favourite. Reading one as
the other transposes the data silently and raises no error. We established which is
which by cross-validation: the political goals data ships in two of these packages and
they agree only when each is read under its own convention.

\begin{thebibliography}{99}

\bibitem{costagomes2022}
Miguel A. Costa-Gomes, Carlos Cueva, Georgios Gerasimou, and Mat\'u\v{s}
Tejiv\v{s}\v{c}\'ak.
Choice, deferral and consistency.
\emph{Quantitative Economics}, 13:1297--1318, 2022.

\bibitem{cotton2021}
Peter Cotton.
Inferring relative ability from winning probability in multientrant contests.
\emph{SIAM Journal on Financial Mathematics}, 12(1), 2021.

\bibitem{goldberg2001}
Ken Goldberg, Theresa Roeder, Dhruv Gupta, and Chris Perkins.
Eigentaste: a constant time collaborative filtering algorithm.
\emph{Information Retrieval}, 4(2):133--151, 2001.

\bibitem{kamishima2003}
Toshihiro Kamishima.
Nantonac collaborative filtering: recommendation based on order responses.
In \emph{Proceedings of KDD}, 2003.

\bibitem{luce1959}
R. Duncan Luce.
\emph{Individual Choice Behavior: A Theoretical Analysis}.
Wiley, 1959.

\bibitem{mattei2013}
Nicholas Mattei and Toby Walsh.
PrefLib: a library of preference data.
In \emph{Algorithmic Decision Theory}, 2013.

\bibitem{mcfaddentrain2000}
Daniel McFadden and Kenneth Train.
Mixed MNL models for discrete response.
\emph{Journal of Applied Econometrics}, 15(5):447--470, 2000.

\bibitem{thurstone1927}
Louis L. Thurstone.
A law of comparative judgment.
\emph{Psychological Review}, 34:273--286, 1927.

\bibitem{yellott1977}
John I. Yellott.
The relationship between Luce's choice axiom, Thurstone's theory of comparative
judgment, and the double exponential distribution.
\emph{Journal of Mathematical Psychology}, 15:109--144, 1977.

\end{thebibliography}
\end{document}
\end{thebibliography}
\end{document}
